\documentclass[aspectratio=169,t,13pt,usenames,dvipsnames]{beamer}
\usetheme[lily]{PaloAlto}
\usecolortheme{dolphin}
\usepackage{mathtools}
\usepackage{tikz}
\usepackage{pgfplots}
\usepackage{tikz-3dplot}
\usepackage{adjustbox}
\usetikzlibrary{quantikz}
\usepackage{xcolor}
\usepackage{bm}
\pgfplotsset{compat=newest}
\logo{%
\stackbox[c][c]{\includegraphics[height=1.15cm]{logo.pdf}\\[-3pt]{\textcopyright~2021}\\{Ron K. Cytron}}}

% \logo{\includegraphics[height=2.5cm]{logo.pdf}}
\def\ColorR#1{\textcolor{BrickRed}{#1}}
\def\ColorG#1{\textcolor{\RCthree}{#1}}
\def\ColorB#1{\textcolor{\RCtwo}{#1}}
\definecolor{links}{HTML}{2E8B57}
\hypersetup{colorlinks,linkcolor=,urlcolor=links}
\def\LinkArrow#1{%
\mbox{$\,$\href{#1}{\mbox{\vrule width 0pt height 0ex depth 0pt$\mapsto$}}}$\,$}
\def\WLOG{WLOG}
\def\True{\mbox{\texttt{True}}}
\def\False{\mbox{\texttt{False}}}
\def\Zero{\mbox{$0$}}
\def\One{\mbox{$1$}}
\def\Not#1{%
\ensuremath{\Overline{\ensuremath{#1}}}}
\def\Xor#1#2{\ensuremath{#1\oplus #2}}
\def\Nand#1#2{%
\ensuremath{\mbox{Nand}(#1,#2)}}
\def\And#1#2{%
\ensuremath{#1 \wedge #2}}
\def\Or#1#2{%
\ensuremath{#1 \vee #2}}
\def\Conj#1{%
\ensuremath{#1^{\star}}}
\def\Star#1{%
\Conj{#1}}
\def\EmptyString{\ensuremath{\lambda}}
\def\Mag#1{%
\ensuremath{|\,#1\,|}}
\def\Prob#1{%
\ensuremath{\Mag{#1}^{2}}}
\def\CNOT#1#2{%
\ensuremath{\NamedGate{CNOT}(#1,#2)}}
\def\CNOTMatrix{%
\begin{pmatrix*}[r]
1 & 0 & 0 & 0 \\
0 & 1 & 0 & 0 \\
0 & 0 & 0 & 1 \\
0 & 0 & 1 & 0\end{pmatrix*}%
}
\def\CCNOT#1#2#3{%
\ensuremath{\mbox{\NamedGate{CNOT}}(#1,#2,#3)}}
\long\def\ToDo#1{\smash{\textcolor{YellowOrange}{#1}}}
\def\ApplyGate#1#2{%
\ensuremath{#1\,#2}}
\def\Inverse#1{\ensuremath{{#1}^{-1}}}
\def\Quote#1{``#1''}
\def\NamedGate#1{%
\mbox{\bfseries{\textrm{#1}}}}
\def\Identity{\NamedGate{I}}
\def\Hadamard{\NamedGate{H}}
\def\PauliX{\NamedGate{X}}
\def\PauliY{\NamedGate{Y}}
\def\PauliZ{\NamedGate{Z}}
\def\SQBG#1#2#3#4#5{%
\ensuremath{%
#1 \begin{pmatrix*}[r] #2 & #3 \\ #4 & #5\end{pmatrix*}}}
\def\HMatrix{%
\SQBG{\RootTwo{}}{1}{1}{1}{-1}}
\def\XMatrix{%
\SQBG{\relax}{0}{1}{1}{0}}
\def\ZMatrix{%
\SQBG{\relax}{1}{0}{0}{-1}}
\def\YMatrix{%
\SQBG{\relax}{0}{-i}{i}{0}}
\def\IMatrix{%
\SQBG{\relax}{1}{0}{0}{1}}
\def\Co#1{\mathbb{C}}
\def\Implies#1#2{\ensuremath{#1 \rightarrow\ #2}}
\def\Domain#1{\ensuremath{\mathcal{#1}}}
\def\CompClass#1{\Domain{#1}}
\def\SpecialX#1{\ensuremath{#1^{\star}}}
\def\Forall#1#2{\ensuremath{\forall\ #1, #2}}
\def\Degrees#1{\ensuremath{#1^{\circ}}}
\def\RootTwo{\ensuremath{\frac{1}{\sqrt{2}}}}
\def\TwoSupOp#1#2#3{%
\ensuremath{\RootTwo{}\left(#1 #3 #2\right)}}
\def\TwoSup#1#2{%
   \TwoSupOp{\ket{#1}}{\ket{#2}}{+}}
\def\QPlus{\TwoSupOp{\QZero}{\QOne}{+}}
\def\QMinus{\TwoSupOp{\QZero}{\QOne}{-}}
\def\Set#1{\ensuremath{\left\{\,#1\,\right\}}}
\def\BigSkip{%

\bigskip

}

\def\MedSkip{%

\medskip

}
\def\SmallSkip{%

\smallskip

}
\def\SineWave#1{%
{%
\pgfmathsetmacro{\Xone}{#1}
\pgfmathsetmacro{\Xtwo}{-(#1)}
\draw (0,0) sin(1,\Xone) cos (2,0) sin(3,\Xtwo) cos (4,0);
}
}
\newcommand{\tolstrut}{%
  \vrule height\dimexpr\fontcharht\font`\A+.1ex\relax width 0pt\relax
}

\DeclareRobustCommand{\Overline}[1]{%
  \ensuremath{\overline{\raisebox{0pt}[1.07\height]{#1}}}%
}
\newenvironment{TIKZP}[1][scale=1.0]{%
\adjustbox{valign=t}\bgroup
\begin{tikzpicture}[#1]
}{%
\end{tikzpicture}
\egroup
}
\long\def\Definition#1{%
\begin{block}{Definition}
#1
\end{block}}
\long\def\Postulate#1#2{%
\begin{block}{Postulate}
\begin{description} \item[#1] #2\end{description}
\end{block}}
\long\def\TwoUnequalColumns#1#2#3#4{%
\begin{columns}%
\begin{column}{#1}#3\end{column}%
\begin{column}{#2}#4\end{column}%
\end{columns}%
}
\long\def\ThreeUnequalColumns#1#2#3#4#5#6{%
\begin{columns}%
\begin{column}{#1}#4\end{column}%
\begin{column}{#2}#5\end{column}%
\begin{column}{#3}#6\end{column}%
\end{columns}%
}
\long\def\ThreeColumns#1#2#3{%
\ThreeUnequalColumns{0.33\textwidth}{0.33\textwidth}{0.33\textwidth}{#1}{#2}{#3}%
}
\long\def\TwoColumns#1#2{%
\TwoUnequalColumns{0.5\textwidth}{0.5\textwidth}{#1}{#2}%
}
\long\def\OnlyRemark#1#2{%
\only<#1>{\Remark{#2}}}

\long\def\Remark#1{%
\begin{block}{Remark}
#1
\end{block}}
\long\def\Example#1{%
\begin{example} #1\end{example}}
\def\VV{\textit{vice versa}}
\def\QZero{\ket{\hbox to 0.7em{\hss 0\hss}}}
\def\QOne{\ket{\hbox to 0.7em{\hss 1\hss}}}
\def\QState#1{\ket{\QName{#1}}}
\def\QName#1{\ensuremath{\psi_{#1}}}
\def\TensProd#1#2{\ensuremath{#1 \bm{\otimes} #2}}
\def\TensSupProd#1#2{\ensuremath{#1^{\bm{\otimes} #2}}}
\def\PZero{\SQB{1}{0}}
\def\POne{\SQB{0}{1}}
\def\PPlus{\ensuremath{\RootTwo{}\SQB{1}{1}}}
\def\PMinus{\ensuremath{\RootTwo{}\SQB{1}{-1}}}
\def\PHad{\ensuremath\frac{1}{\sqrt{2}}\SQB{1}{1}}
\def\PHadm{\ensuremath\frac{1}{\sqrt{2}}\SQB{1}{-1}}
\def\Intersect#1#2{\ensuremath{#1 \cap #2}}
\def\Union#1#2{\ensuremath{#1 \cup #2}}
\def\Set#1{\ensuremath{\left\{\,#1\,\right\}}}
\def\AllBits#1#2{\ensuremath{#1\in \Set{0,1}^{#2}}}
\def\DQB#1#2#3#4{%
\ensuremath{\begin{pmatrix*}[r] #1 \\ #2 \\ #3 \\ #4\end{pmatrix*}}}
\def\QQB#1#2#3#4#5#6#7#8{%
\ensuremath{\begin{pmatrix*}[r] #1 \\ #2 \\ #3 \\ #4 \\ #5 \\ #6 \\ #7 \\ #8\end{pmatrix*}}}
\def\SQB#1#2{%
\ensuremath{\begin{pmatrix*}[r] #1 \\ #2\end{pmatrix*}}}
\def\CQB#1#2{%
\ensuremath{\begin{pmatrix*}[r] #1 #2\end{pmatrix*}}}
\def\CDQB#1#2#3#4{%
\ensuremath{\begin{pmatrix*}[r] #1 #2 #3 #4\end{pmatrix*}}}
\def\TZPointDiam{2pt}
\def\TZPoint#1#2#3{%
\draw[fill=black] (#1) circle (\TZPointDiam) node[#3] {#2} ;}
\def\UnitComplexCircle{%
\draw [<->] (-1.5, 0) -- (1.5, 0)  node[right] {$\Re{}$} ;
   \draw [<->] (0,-1.5) -- (0, 1.5) node[above] {$\Im{}$} ;

   \draw (0, 0) circle (1) ;
   }
   \def\TZPEast{\TZPoint{1,0}{$+1$}{above right}}
   \def\TZPNorth{\TZPoint{0,1}{$i$}{above right}}
   \def\TZPWest{\TZPoint{-1,0}{$-1$}{above left}}
   \def\TZPSouth{\TZPoint{0,-1}{$-i$}{below right}}
\def\TZText#1#2#3{%
  \draw (#1) node [#3] {#2};
}
\def\Qiskit{\texttt{qiskit}}
\def\TwoToThe#1{\ensuremath{2^{#1}}}
\def\Exp#1{\ensuremath{e^{#1}}}
\def\ExpPhase#1{\Exp{i\,#1}}
\def\ExpNegPhase#1{\Exp{-i\,#1}}
%%
%% Not sure this is used
%%
\def\PhaseLabel#1{\phase{\fbox{#1}}}
%%
\def\ColHead#1{\begin{center}\bf #1\end{center}}
\def\Complex#1#2{\ComplexGen{#1}{#2}{+}}
\def\ComplexDiff#1#2{\ComplexGen{#1}{#2}{-}}
\def\ImaginaryM#1{\ComplexDiff{\relax}{#1}}
\def\ComplexGen#1#2#3{%
\ensuremath{#1 #3 i\relax #2}}
\def\Polar#1#2{\ensuremath{#1\hbox to 0.1pt{\hss}\Exp{i\relax#2}}}
%%
%% #1 -- lower left x
%% #2 -- lower left y
%% #3 -- upper right x
%% #4 -- upper right y
%% #5 -- num lines
\def\PFilter#1#2#3#4#5{

    \draw (#1,#2) -- (#3,#2) -- (#3,#4) -- (#1,#4) -- cycle;
    \foreach \i in {1,...,#5}
    {
        \pgfmathsetmacro{\PFilterInc}{#1 + (#3-#1) / #5 * \i}
        % \edef\PFilterInc{\pgfmathresult}
        \draw (\PFilterInc, #2) -- (\PFilterInc, #4);
    }
}
%% #1 - x
%% #2 - y
%% #3 - degrees
%% #4 - options
%% #5 - radius
\def\RadiantArrowsR#1#2#3{%
\foreach \p in {0, #1,...,360} {
   \draw[#2] (0,0)  -- (\p:#3) ;
}
}
\def\RadiantArrows#1#2{%
   \RadiantArrowsR{#1}{#2}{1}
}
\newcommand\RotateTheta[1][scale=1.0]{%
\SquareOutline{}
\begin{scope}[#1]
\fill[white] (0.5,0.5) circle (0.5);
\draw (0.5,0.5) circle (0.5);
\draw[->] (1.0,0.5) arc(0:45:0.5);
\draw (0.5,0.5) node {$\theta$};
\end{scope}
}
\def\SquareOutline{%
\path (0,0) rectangle (1,1);%
}%
\newcommand\Qif[1][scale=1.0]{%
\begin{scope}[#1]
\SquareOutline{}
\clip (0,0) rectangle(1,1);
\RotateAroundCenter{45}{%
\begin{scope}
\draw (0.15,0.15) rectangle(0.85,0.85);
\end{scope}}
\end{scope}
}
\newcommand\LightSource[1][scale=1.0]{%
\begin{scope}[#1]
\SquareOutline{}
\PhotonGun{}
\begin{scope}[shift={(0.75,0.5)}]
    \RadiantArrowsR{22.5}{color=red}{0.25}
\end{scope}
\end{scope}
}
\def\PhotonGun{%
\fill (0,0.375) rectangle (0.5, 0.625);
}
\def\HPolar{%
\begin{scope}[color=red]
\SquareOutline{}
\draw[thick] (0.5,0.4) -- ++(0.35,0)
             ++(-0.35,0.1) -- ++(0.35,0)
             ++(-0.35,0.1) -- ++(0.35,0);
% \fill (0.90,0.5) circle (0.05);
\end{scope}
}
\def\VPolar{%
\begin{scope}[color=red]
\draw[thick] (0.6,0.375) -- ++(0,.250)
             ++(0.1,-.250) -- ++(0,.250)
             ++(0.1,-.250) -- ++(0,.250);
% \fill (0.95,0.5) circle (0.05);
\end{scope}
}
\newcommand\HPolarizedLightSource[1][scale=1.0]{%
\begin{scope}[#1]
\SquareOutline{}
\PhotonGun{}
\HPolar{}
\end{scope}
}
\newcommand\VPolarizedLightSource[1][scale=1.0]{%
\begin{scope}[#1]
\SquareOutline{}
\PhotonGun{}
\VPolar{}
\end{scope}
}
\newcommand\EVBomb[1][scale=1.0]{%
\begin{scope}[#1]
\fill (0.5,0.5) circle(0.5);
\end{scope}%
}
\newcommand\EVBoom[1][scale=1.0]{%
\begin{scope}[#1]
\fill[white] (0.5,0.5) circle(0.5);
\draw[red,thick] (0.5,0.5) circle(0.5);
\clip (0,0) rectangle(1,1);
\begin{scope}
\node[scale=0.35,xshift=28pt,yshift=28pt,starburst,draw,red,fill=orange,line width=1.5pt]
{BOOM!};
\end{scope}
\end{scope}
}
\newcommand\Measurement[1][scale=1.0]{%
\begin{scope}[#1]
\draw[color=white] (0,0) rectangle(1,1);
\begin{scope}
\clip (0.25,0.25) rectangle (1,1);
\draw (0,0) circle(0.9);
\draw (0,0) circle(1.0);
\end{scope}
\draw[thick] (0,0) --(0,1);
\draw[->] (0.25,0.25) -- (0.9,0.9);
\end{scope}
}
\def\TRectangle#1#2#3#4#5#6{%
\begin{scope}[#6]
\draw [fill=#5] (#1,#2) rectangle (#3,#4);
\end{scope}
}
%% #1 -- llx
%% #2 -- lly
%% #3 -- urx
%% #4 -- ury
%% #5 -- angle
\newcommand\Mirror[1][rotate=0]{%
\clip (0,0) rectangle(1,1);
\begin{scope}[#1]
   \SquareOutline{}
   \TRectangle{0}{.4375}{1}{.5625}{Gray}{rotate=0}
\end{scope}
}
\newcommand\BeamSplitter[1][scale=1.0]{%
\begin{scope}[#1]
    \SquareOutline{}
    \TRectangle{0}{.375}{1}{.625}{SkyBlue}{rotate=0}
\end{scope}
}
\def\MZ{%
\LightSource{}
\Shift{2}{0}{\RotateAroundCenter{-45}{\BeamSplitter{}}}
\Shift{2}{-2}{\RotateAroundCenter{-45}{\Mirror{}}}
\Shift{6}{0}{\RotateAroundCenter{-45}{\Mirror{}}}
\Shift{6}{-2}{\RotateAroundCenter{-45}{\BeamSplitter{}}}
\Shift{8}{-2}{\Measurement[color=orange]{}}
\Shift{6}{-4}{\RotateAroundCenter{-90}{\Measurement[color=\RCtwo]{}}}
}
\def\Shift#1#2#3{%
\begin{scope}[shift={(#1,#2)}] #3\end{scope}
}
\def\RotateAroundCenter#1#2{%
\begin{scope}[rotate around={#1:(0.5,0.5)}]
  #2
\end{scope}
}
\def\Vskip#1{\mbox{}\vskip #1\mbox{}}
\def\Hskip#1{\mbox{}\hskip #1\mbox{}}
\author{Ron K.~Cytron}
\institute{Washington University\\Saint Louis, Missouri\\[2em] \textcopyright~2021 All rights reserved by the author}
\setbeamertemplate{footline}[frame number]{}
%\setbeamertemplate{footline}[text line]{%
%  \parbox{\linewidth}{\vspace*{-8pt}\textcopyright~2021\hfill\insertshortauthor\hfill\insertpagenumber}}
%% #1 -- lecture number
%% #2 -- lecture title
%% #3 -- lecture subtitle
%% #4 -- lecture label
\def\SetTitle#1#2#3#4{%
   \title{#2}%
   \subtitle{#3}%
   \date{\today}%
   %\begin{frame}\maketitle\end{frame}%
   \frame{\titlepage}%
   \lecture[#1]{#2: #3}{lec:#4}%
}
\def\Kaye{\href{https://dl.acm.org/doi/10.5555/1206629}{Kaye}}
\def\MikeIke{\href{https://dl.acm.org/doi/10.5555/1972505}{Nielson~and~Chuang}}

\newcounter{ProtocolDialogStep}
\newenvironment{ProtocolDialog}[3]{%
\def\Incr{\stepcounter{ProtocolDialogStep}}
\def\Ref##1{\Incr{}\ThreeUnequalColumns{#1}{#2}{#3}{\arabic{ProtocolDialogStep}. ##1}{\relax}{\relax}}%
\def\Alice##1{\Incr{}\ThreeUnequalColumns{#1}{#2}{#3}{\relax}{\arabic{ProtocolDialogStep}. ##1}{\relax}}%
\def\All##1##2##3{%
\Incr{}
\ThreeUnequalColumns{#1}{#2}{#3}{\arabic{ProtocolDialogStep}. ##1}{##2}{##3}
}
\def\Bob##1{\Incr{}\ThreeUnequalColumns{#1}{#2}{#3}{\relax}{\relax}{
\arabic{ProtocolDialogStep}. ##1}}%
\setcounter{ProtocolDialogStep}{0}%
\ThreeUnequalColumns{#1}{#2}{#3}{\textbf{Ref}}{\textbf{Alice}}{\textbf{Bob}}
}{%
}
%% #1 -- options
%% #2 -- width
%% #3 -- height
%% #4 -- signals
\newenvironment{GateBox}[4][scale=1.0]{%
\edef\Signals{#4}
\edef\Height{#3}
\edef\Width{#2}
\pgfmathsetmacro{\Max}{\Signals}
\pgfmathsetmacro{\Vsep}{\Height / \Max}
\pgfmathsetmacro{\Wlen}{\Width / 5}
\def\ComputeMyY##1{%
   \pgfmathsetmacro{\MyY}{\Height - 0.5 * \Vsep - ##1*\Vsep}}
\def\Input##1##2{%
   \pgfmathsetmacro{\Xone}{-\Wlen}
   \pgfmathsetmacro{\Xtwo}{0}
   \pgfmathsetmacro{\MyY}{\Height - 0.5 * \Vsep - ##1*\Vsep}
   \draw[->] (\Xone,\MyY) node[left] {##2} -- (\Xtwo,\MyY);
}
\def\Output##1##2{%
   \pgfmathsetmacro{\Xone}{\Width}
   \pgfmathsetmacro{\Xtwo}{\Width + \Wlen}
   \pgfmathsetmacro{\MyY}{\Height - 0.5 * \Vsep - ##1*\Vsep}
   \draw[->] (\Xone,\MyY) -- (\Xtwo,\MyY) node[right] {##2};
}
\def\BoxLabel##1{%
\node[draw, fit={(0,0) (\Width,\Height)}, inner sep=0pt, label=center:\mbox{##1}] (A) {};
}
\begin{TIKZP}[#1]
   \draw (0,0) rectangle(#2,#3);
}{%
\end{TIKZP}
}
\def\PhiTheta#1#2{%
\mbox{\ensuremath{\textcolor{\RCone}{\phi = #1}, \textcolor{\RCtwo}{\theta = #2}}}}
%
% #2 -- orcacle name
% #3 -- upper left input
% #4 -- lower left input
% #5 -- upper right output
% #6 -- lower right output
%
\newenvironment{Oracle}[6][scale=1.0]{%
\begin{GateBox}[#1]{2.5}{1.5}{2}
\BoxLabel{#2}
\Input{0}{#3}
\Input{1}{#4}
\Output{0}{#5}
\Output{1}{#6}
}{%
\end{GateBox}
}
\newenvironment{TTable}[3]{%
\def\TRow##1##2##3{%
##1 & ##2 & ##3 \\}
\begin{tabular}{c|cc}
\TRow{#1}{#2}{#3}\hline
}{%
\end{tabular}
}
%%
%%  1 name
%%  2 one x
%%  3 one f(x)
%%  4 another x
%%  5 another f(x)
\def\SBitTable#1#2#3#4#5{%
\begin{center}
    \begin{TTable}{#1}{$x$}{$f(x)$}
    \TRow{ }{#2}{#3}
    \TRow{ }{#4}{#5}
    \end{TTable}
\end{center}}
%%
%%
\newenvironment{UTable}[6]{%
\def\URow##1##2##3##4##5##6{%
##1 & ##2 & ##3  & ##4 & ##5 & ##6\\}
\begin{tabular}{c|cc|c|cc}
& \multicolumn{2}{|c|}{Inputs} &
& \multicolumn{2}{|c}{Outputs} \\
\URow{#1}{#2}{#3}{#4}{#5}{#6}\hline
}{%
\end{tabular}
}
%%
%%  1 name
%%  2 one x
%%  3 one f(x)
%%  4 another x
%%  5 another f(x)
\def\SBitTable#1#2#3#4#5{%
\begin{center}
    \begin{TTable}{#1}{$x$}{$f(x)$}
    \TRow{ }{#2}{#3}
    \TRow{ }{#4}{#5}
    \end{TTable}
\end{center}}

\newenvironment{BlochSphere}[3][scale=1.0]{
    \bgroup%
    \def\BlochSpherePhiTheta{%
    \tdplotsetrotatedcoords{0}{0}{0};O
        \path[tdplot_rotated_coords] (1.0,0,0) arc (0:115:1.0) node(phi) {};
        \draw[tdplot_rotated_coords,thick,\RCone] (0.3,0,0) arc (0:115:0.3);
        \draw[tdplot_rotated_coords,\RCone,thick,dashed] (0,0,0) -- (phi);
        \draw[\RCone] (0.4,0.4,0) node {$\phi$};
        
        \tdplotsetrotatedcoords{25}{-90}{0};
        \path[tdplot_rotated_coords] (1,0,0) arc (0:45:1) node(theta) {};
        \draw[tdplot_rotated_coords,thick,\RCtwo] (0.4,0,0) arc (0:45:0.4);
        \draw[tdplot_rotated_coords,\RCtwo,thick] (0,0,0) -- (theta);
        \draw[\RCtwo] (-0.4,0,0.3) node {$\theta$};
        
        
    }
    \def\TZPointDiam{0.7pt}%
    \def\ShowAxes{%
    \draw[dashed, tdplot_rotated_coords, black] (1,0,0) node {{\Large\bf x}} (0,1,0) node {{\Large\bf y}} (0,0,1) node {{\Large\bf z}};
    }%
    \tdplotsetmaincoords{#2}{#3}%
    \begin{tikzpicture}[tdplot_main_coords, #1]
        
        
        \shade[ball color = lightgray, opacity = 0.5] (0,0,0) circle (1cm);
        
        %%
        %% The equator
        %%
        \tdplotsetrotatedcoords{0}{0}{0};
        \draw[dashed, tdplot_rotated_coords, gray] (0,0,0) circle (1);
        
        %%
        %% Arc in the x-z plane
        %%
        \tdplotsetrotatedcoords{90}{90}{90};
        \draw[dashed, tdplot_rotated_coords, gray] (1,0,0) arc (0:180:1);
        
        %%
        %% Arc in the y-z plane
        %%
        \tdplotsetrotatedcoords{0}{90}{90};
        \draw[dashed, tdplot_rotated_coords, gray] (1,0,0) arc (0:180:1);
        
        %\tdplotsetrotatedcoords{90}{90}{90};
       % \draw[dashed, tdplot_rotated_coords, black] (1,0,0) node {{\Large\bf x}} (0,1,0) node {{\Large\bf y}} (0,0,1) node {{\Large\bf z}}; 
        
        \draw[dashed, gray] (0,0,0) -- (-1,0,0);
        \draw[dashed, gray] (0,0,0) -- (0,-1,0);
        \draw[dashed, gray] (0,0,0) -- (0,0,-1);
        
        \draw[-stealth] (0,0,0) -- (1.30,0,0);
        \draw (1.45,0,0) node {$x$};
        \draw[-stealth] (0,0,0) -- (0,1.30,0);
        \draw(0,1.45,0) node {$y$};
        \draw[-stealth] (0,0,0) -- (0,0,1.30);
        \draw (0,0,1.45) node {$z$};
        
}{%
    \end{tikzpicture}
    \tdplotsetmaincoords{0}{0}
    \egroup
}

\long\def\Exercise#1#2{%
\relax}
\def\QED{\begin{flushright}\ensuremath{\square}\end{flushright}}
\def\UGateA{%
\ensuremath{\cos(\theta/2)}}
\def\UGateB{%
\ensuremath{-\ExpPhase{\lambda}\sin(\theta/2)}}
\def\UGateC{%
\ensuremath{\ExpPhase{\phi}\sin(\theta/2)}}
\def\UGateD{%
\ensuremath{\ExpPhase{(\phi+\lambda)}\cos(\theta/2)}}
\def\UGate{%
\SQBG{\relax}{\UGateA{}}{\UGateB{}}{\UGateC{}}{\UGateD{}}}%
\def\FCirc#1{%
\begin{tikzpicture}[scale=#1]%
\SquareOutline{}%
\fill (0.5,0.5) circle(0.4);%
\end{tikzpicture}%
}%
\def\SCirc{\FCirc{0.3}}%
\newenvironment{BBKey}[0]{%
\def\STD{%
\begin{tikzpicture}[scale=0.35]%
\SquareOutline{}%
\draw[ultra thick,purple] (0,0.5) -- (1,0.5);%
\draw[ultra thick,purple] (0.5,0) -- (0.5,1);%
\end{tikzpicture}}%
\def\HDM{%
\begin{tikzpicture}[scale=0.30]%
\SquareOutline{}%
\draw[ultra thick,purple] (0,0) -- (1,1);%
\draw[ultra thick,purple] (0,1) -- (1,0);%
\end{tikzpicture}}%
\def\BBUp{%
\begin{tikzpicture}[scale=0.30]%
\SquareOutline{}%
\draw[->,thick,purple] (0.5,0) -- (0.5,1);%
\end{tikzpicture}}%
\def\BBRt{%
\begin{tikzpicture}[scale=0.30]%
\SquareOutline{}%
\draw[->,thick,purple] (0,0.5) -- (1,0.5);%
\end{tikzpicture}}%
\def\BBNe{%
\begin{tikzpicture}[scale=0.30]%
\SquareOutline{}%
\draw[->,thick,purple] (0,0) -- (1,1);%
\end{tikzpicture}}%
\def\BBSe{%
\begin{tikzpicture}[scale=0.30]%
\SquareOutline{}%
\draw[->,thick,purple] (0,1) -- (1,0);%
\end{tikzpicture}}%
\def\BBBasis{%
\begin{tabular}{c|c|c}
Basis & 0 & 1 \\
\hline
\STD{} & \BBUp{} & \BBRt{} \\
\HDM{} & \BBNe{} & \BBSe{}
\end{tabular}%
}%
}{%
\relax%
}
\def\CFCirc#1{%
  \textcolor{#1}{\FCirc{0.30}}}
\def\Agree{%
  \textcolor{green}{\FCirc{0.30}}}%
\def\DisAgree{%
   \textcolor{red}{\FCirc{0.30}}}%
\def\MetaAgree{%
   \textcolor{yellow}{\FCirc{0.30}}}%
\newenvironment{EKey}[0]{%
\def\ZZ{%
\ZRot{0}%
}%
\def\ZZPO##1{%
\ZRot{\frac{\pi}{##1}}%
}%
\def\ZZNPO##1{%
\ZRot{\ensuremath{\,-\frac{\pi}{##1}}}%
}%
\def\ZRot##1{%
\ensuremath{Z_{##1}}%
}%
}{%
}
\def\AltAlert#1#2#3{%
\alt<#1>{\alert{#2}}{#3}}
\def\RCone{Rhodamine}
\def\RCtwo{CornflowerBlue}
\def\RCthree{LimeGreen}
\def\RCfour{Tan}
\def\ColorOne#1{\textcolor{\RCone}{#1}}
\def\ColorTwo#1{\textcolor{\RCtwo}{#1}}
\def\ColorThree#1{\textcolor{\RCthree}{#1}}
\def\ColorFour#1{\textcolor{\RCfour}{#1}}
\def\NegF#1{\ensuremath{(-1)^{f(#1)}}}
\def\RootTwoN#1{\ensuremath{\frac{1}{\sqrt{2^#1}}}}



% \AtBeginLecture{\frame{\Large \insertlecture}}
%%
%% test
%%
\begin{document}
\SetTitle{21}{Deutsch's Problem}{Is $f(x)$ constant or balanced?}{20}

\begin{frame}{Deutsch's problem}{Formulation}

\begin{itemize}
    \item We are given an oracle $f(x): \Set{0,1} \mapsto \Set{0,1}$
    \item We can characterize $f(x)$ as follows:
    \begin{description}
        \item[constant]   \Forall{x}{f(x)=k}, where $k$ is $0$ or $1$
        \item[balanced]   For half of its inputs, $f(x)=1$ and for the other half $f(x)=0$
    \end{description}
    \item Given the domain and range of $f(x)$, the function must be one of the above.
\end{itemize}
    
\end{frame}

\begin{frame}{Possible functions}{There are only four.}

\TwoColumns{%
\visible<2->{\begin{center}\alert{Constant}\end{center}}
\SBitTable{A}{0}{0}{1}{0}
\SBitTable{B}{0}{1}{1}{1}
}{%
\visible<3>{\begin{center}\alert{Balanced}\end{center}}
\SBitTable{C}{0}{0}{1}{1}
\SBitTable{D}{0}{1}{1}{0}
}
    
\end{frame}

\begin{frame}{Form of the oracle}{It has to be reversible.}
\TwoUnequalColumns{0.6\textwidth}{0.4\textwidth}{%
\begin{itemize}[<+->]
    \item We consider a quantum gate $U_f$ that implements $f(x)$.
    \item It accepts $x$ and an \href{https://en.wikipedia.org/wiki/Ancilla_bit}{ancillary} input $y$.
    \item It copies $x$ onto the top output qubit.
    \item It computes \alert{\Xor{y}{f(x)}} on the bottom output qubit.
    \item When $y=\QZero{}$ the bottom qubit is $\Xor{0}{f(x)}=f(x)$.
    \item When $y=\QOne{}$ the bottom qubit is $\Xor{1}{f(x)}=\Not{f(x)}$.
\end{itemize}
}{%
\BigSkip{}
\only<1>{%
\begin{Oracle}[scale=0.75]{$U_f$}{\relax}{\relax}{\relax}{\relax}
\end{Oracle}}%
\only<2>{%
\begin{Oracle}[scale=0.75]{$U_f$}{$x$}{$y$}{\relax}{\relax}
\end{Oracle}}%
\only<3>{%
\begin{Oracle}[scale=0.75]{$U_f$}{$x$}{$y$}{$x$}{\relax}
\end{Oracle}}%
\only<4>{%
\begin{Oracle}[scale=0.75]{$U_f$}{$x$}{$y$}{$x$}{\alert{\Xor{y}{f(x)}}}
\end{Oracle}}%
\only<5>{%
\begin{Oracle}[scale=0.75]{$U_f$}{$x$}{\QZero{}}{$x$}{$f(x)$}
\end{Oracle}}%
\only<6>{%
\begin{Oracle}[scale=0.75]{$U_f$}{$x$}{\QOne{}}{$x$}{\Not{$f(x)$}}
\end{Oracle}}%
}    
\end{frame}



\begin{frame}{Realizing the oracle for \Quote{A}: $f(x)=0$}{We design a circuit that meets the classical specification.}

\TwoColumns{%
\Vskip{-3em}\begin{center}
\begin{TTable}{A}{$x$}{$f(x)$}
\TRow{ }{0}{0}
\TRow{ }{1}{0}
\end{TTable}
\end{center}
\BigSkip{}
\visible<2->{%
\begin{center}
\begin{UTable}{$U_A$}{$x$}{$y$}{$f(x)$}{$x$}{\Xor{y}{f(x)}}
\URow{}{0}{0}{0}{0}{0}
\URow{}{0}{1}{0}{0}{1}
\URow{}{1}{0}{0}{1}{0}
\URow{}{1}{1}{0}{1}{1}
\end{UTable}\end{center}}
}{%
\begin{Oracle}[scale=0.75]{$U_A$}{$x$}{$y$}{$x$}{\Xor{y}{f(x)}}
\end{Oracle}

\BigSkip{}
\visible<3->{%
\Forall{x}{f(x)=0}, so $\Xor{y}{f(x)}=\Xor{y}{0}=y$.  
\MedSkip{}
Thus, $U_A$ is just the identity function, so our quantum circuit doesn't have to do anything to its inputs to realize $U_A$.}

\BigSkip{}
\visible<4>{%
\begin{center}
\adjustbox{valign=t}{\begin{quantikz}
\lstick{\ket{x}} &  \qw & \qw \rstick{\ket{x}} \\
\lstick{\ket{y}} &  \qw & \qw \rstick{\ket{\Xor{y}{f(x)}}}
\end{quantikz}}\end{center}
}
}
    
\end{frame}

\begin{frame}{Realizing the oracle for \Quote{B}: $f(x)=1$}{We design a circuit that meets the classical specification.}

\TwoColumns{%
\Vskip{-3em}\begin{center}
\begin{TTable}{B}{$x$}{$f(x)$}
\TRow{ }{0}{1}
\TRow{ }{1}{1}
\end{TTable}
\end{center}
\BigSkip{}
\visible<2->{%
\begin{center}
\begin{UTable}{$U_B$}{$x$}{$y$}{$f(x)$}{$x$}{\Xor{y}{f(x)}}
\URow{}{0}{0}{1}{0}{1}
\URow{}{0}{1}{1}{0}{0}
\URow{}{1}{0}{1}{1}{1}
\URow{}{1}{1}{1}{1}{0}
\end{UTable}\end{center}}
}{%
\begin{Oracle}[scale=0.75]{$U_B$}{$x$}{$y$}{$x$}{\Xor{y}{f(x)}}
\end{Oracle}

\BigSkip{}
\visible<3->{%
\Forall{x}{f(x)=1}, so $\Xor{y}{f(x)}=\Xor{y}{1}=\Not{y}$.  
\MedSkip{}
Thus, $U_B$ is realized by complementing $y$.}

\BigSkip{}
\visible<4>{%
\begin{center}
\adjustbox{valign=t}{\begin{quantikz}
\lstick{\ket{x}} &  \qw & \qw \rstick{\ket{x}} \\
\lstick{\ket{y}} &  \gate{X} & \qw \rstick{\ket{\Xor{y}{f(x)}}}
\end{quantikz}}\end{center}
}
}
    
\end{frame}

\begin{frame}{Realizing the oracle for \Quote{C}: $f(x)=x$}{We design a circuit that meets the classical specification.}

\TwoColumns{%
\Vskip{-3em}\begin{center}
\begin{TTable}{C}{$x$}{$f(x)$}
\TRow{ }{0}{0}
\TRow{ }{1}{1}
\end{TTable}
\end{center}
\BigSkip{}
\visible<2->{%
\begin{center}
\begin{UTable}{$U_C$}{$x$}{$y$}{$f(x)$}{$x$}{\Xor{y}{f(x)}}
\URow{}{0}{0}{0}{0}{0}
\URow{}{0}{1}{0}{0}{1}
\URow{}{1}{0}{1}{1}{1}
\URow{}{1}{1}{1}{1}{0}
\end{UTable}\end{center}}
}{%
\begin{Oracle}[scale=0.75]{$U_C$}{$x$}{$y$}{$x$}{\Xor{y}{f(x)}}
\end{Oracle}

\BigSkip{}
\visible<3->{%
\Forall{x}{f(x)=x}, so $\Xor{y}{f(x)}=\Xor{y}{x}$.}%

\BigSkip{}
\only<3>{%
The matrix may look familiar?
\[U_{C} = \CNOTMatrix{} \]
}%
\only<4>{%
Thus, a \NamedGate{CNOT} gate will flip $y$ when $x$ is~$1$:
\begin{center}
\adjustbox{valign=t}{\begin{quantikz}
\lstick{\ket{x}} &  \ctrl{1} & \qw \rstick{\ket{x}} \\
\lstick{\ket{y}} &  \targ{} & \qw \rstick{\ket{\Xor{y}{f(x)}}}
\end{quantikz}}\end{center}
}
}
    
\end{frame}

\begin{frame}{Realizing the oracle for \Quote{D}: $f(x)=\Not{x}$}{We design a circuit that meets the classical specification.}

\TwoColumns{%
\Vskip{-3em}\begin{center}
\begin{TTable}{D}{$x$}{$f(x)$}
\TRow{ }{0}{1}
\TRow{ }{1}{0}
\end{TTable}
\end{center}
\BigSkip{}
\visible<2->{%
\begin{center}
\begin{UTable}{$U_D$}{$x$}{$y$}{$f(x)$}{$x$}{\Xor{y}{f(x)}}
\URow{}{0}{0}{1}{0}{1}
\URow{}{0}{1}{1}{0}{0}
\URow{}{1}{0}{0}{1}{0}
\URow{}{1}{1}{0}{1}{1}
\end{UTable}\end{center}}
}{%
\begin{Oracle}[scale=0.75]{$U_D$}{$x$}{$y$}{$x$}{\Xor{y}{f(x)}}
\end{Oracle}

\BigSkip{}
\visible<3->{%
\Forall{x}{f(x)=\Not{x}}, so $\Xor{y}{f(x)}=\Xor{y}{\Not{x}}=\Not{\Xor{y}{x}}$.}%

\BigSkip{}
\only<3>{%
This just like the previous case, except the output is the complement of what we had before.
}%
\only<4>{%
Thus, we add an \PauliX{} gate after the \NamedGate{CNOT} gate:
\begin{center}
\adjustbox{valign=t}{\begin{quantikz}
\lstick{\ket{x}} &  \ctrl{1} & \qw & \qw \rstick{\ket{x}} \\
\lstick{\ket{y}} &  \targ{} & \gate{X} & \qw \rstick{\ket{\Xor{y}{f(x)}}}
\end{quantikz}}\end{center}
}
}
    
\end{frame}

\begin{frame}{Is $f(x)$ constant or balanced?}{Classical approach}

\TwoColumns{%
\begin{itemize}[<+->]
    \item If we perform one evaluation of $f$, say on $0$, then we do not know if $f(x)$ is always $f(0)$ or if $f(1)$ may be different.
    \item Classically it takes \emph{two} queries of $f$ to determine the answer to this question.
    \visible<3-4>{Here $f(x)$ is \alt<3>{constant}{balanced}.}
\end{itemize}
}{%
\begin{align*}
    f(0) &= 0 \\
    \visible<2->{f(1) &= \alt<4>{1}{\alt<3>{0}{\ ?}}}
\end{align*}
}
\BigSkip{}
\visible<5>{We will see that a quantum circuit for $f$ can reveal whether the function is constant or balanced with \emph{a single} query!}
\end{frame}

\begin{frame}{Using superposition}{We can arrange for the quantum circuit to produce all possible answers.}

\Vskip{-3em}\begin{center}
\begin{GateBox}[scale=1.0]{2}{1.2}{2}
\BoxLabel{$U_f$}
\Input{0}{\alt<1>{\QPlus{}}{\ket{+}}}
\Input{1}{\alt<16->{\alert{\ket{-}}}{\QZero{}}}
\ComputeMyY{0}\draw (0,\MyY) node[right] {$x$};
\ComputeMyY{1}\draw (0,\MyY) node[right] {$y$};
\Output{0}{\alert<12-14>{$x$}}
\Output{1}{\alert<12-14>{\Xor{y}{f(x)}}}
\end{GateBox}
\end{center}
\TwoColumns{%
\only<1-3>{%
\begin{itemize}
    \item<1-> The input is the superposition of the two possible basis states.
    \item<2-> We know this state as \ket{+}.
    \item<3-> Recall we can create $\ket{+}=\QPlus{}$ using a Hadamard gate.
\end{itemize}}
}{%
\only<3>{%
\begin{center}
\adjustbox{valign=t}{\begin{quantikz}
\lstick{\QZero{}} &  \gate{H} &  \qw \rstick{\ket{+}}
\end{quantikz}}%
\end{center}}%
}%

\only<4-6>{%
    $U_f$ produces
        $\TensProd{x}{(\Xor{y}{f(x)})}  =
        \TensProd{x}{(\Xor{0}{f(x)})} 
        = \TensProd{x}{f(x)}
        = \alert{\ket{x\,f(x)}}$

\SmallSkip{}
By linearity, with $x$ in the superposition \QPlus{} and $y=\QZero{}$
    \begin{align*}
    U_f(\TensProd{\QPlus}{0}) &= \RootTwo{}U_f(\ket{00}+\ket{10}) \\
    \visible<5->{ &= \RootTwo{}\left(
        U_f(\ket{00}) + U_f(\ket{10})\right) \\}
       \visible<6->{ &= \frac{\ket{0\,f(0)} +
        \ket{1\,f(1)}}{\sqrt{2}}}
    \end{align*}}%
\only<7->{%
\visible<7-16>{
\[ U_{f}(\ket{+\,0}) = \TwoSup{0\,\textcolor<14>{\RCone}{f(0)}}{1\,\textcolor<14>{\RCtwo}{f(1)}} \]}}
\only<7-11>{%
\Vskip{-2em}\begin{itemize}
    \item<7-> We forced $U_f$ to produce \emph{both} outputs.
    \item<8-> But what happens if we measure?
    \item<9-> We randomly and uniformly see \emph{either} \ket{0\,f(0)} or \ket{1\,f(1)}.
    \item<10-> It could take \emph{more} than two evaluations to see both answers!
    \item<11-> Even if we measure only the top qubit, the bottom qubit collapses too.
\end{itemize}
}%
\only<12-14>{%
\Vskip{-2em}\begin{itemize}
    \item<12-> It is misleading to see the \alert<12>{output} of $U_f$ labeled to look like a tensor product.  $U_f$ will generate the outputs as specified, but the result is not necessarily expressable as \TensProd{x}{(\Xor{y}{f(x)})}.
    \item<13-> Without analysis, we cannot know if the output is tensor-factorable.
    \item<14-> Consider \textcolor{\RCone}{$f(0)=0$} and \textcolor{\RCtwo}{$f(1)=1$}.  Then $U_f$ creates \alert<14>{\TwoSup{00}{11}} as the (entangled) output, which we know cannot be tensor factored!
\end{itemize}
}%
\only<15->{%
\begin{itemize}
    \item<15-> While the information is in the output, we cannot recover it by measurement.
    \item<16-> We need to try something different, using \alert{interference} to get the result we need.
\end{itemize}
}
\end{frame}

\begin{frame}{Creating \ket{-}}{A few approaches}

\TwoUnequalColumns{0.6\textwidth}{0.4\textwidth}{%
\Vskip{-3em}\begin{itemize}[<+->]
    \item Recall 
    \begin{itemize}
        \item $\ket{+}=\QPlus{}$
        \item $\ket{-}=\QMinus{}$
    \end{itemize}
    \item Use \PauliZ{} on \ket{+}
    \[
    \ZMatrix{}\PPlus{} = \PMinus{}=\ket{-}\]
    \item Use \Hadamard{} on \QOne{}
    \[ \HMatrix{}\POne{} = \PMinus{}=\ket{-}
    \]
\end{itemize}
}{%
\BigSkip{}\BigSkip{}
\visible<4>{%
\begin{center}
\adjustbox{valign=t}{\begin{quantikz}
\lstick{\QZero{}} &  \gate{H} & \gate{Z}& \qw \rstick{\ket{-}}
\end{quantikz}}\end{center}}%
\BigSkip{}
\BigSkip{}
\BigSkip{}
\visible<5>{%
\begin{center}
\adjustbox{valign=t}{\begin{quantikz}
\lstick{\QOne{}} &  \gate{H} & \qw & \qw \rstick{\ket{-}}
\end{quantikz}}\end{center}}%
}
    
\end{frame}



\begin{frame}{Deutsch's algorithm}{Uses interference to find constant vs. balanced functions}

\Vskip{-4em}\TwoUnequalColumns{0.5\textwidth}{0.5\textwidth}{%
\begin{center}
\Vskip{-3em}\adjustbox{valign=t, width=\textwidth}{\begin{quantikz}
\lstick{\QZero{}} & \qw\slice{\alert<2>{\QState{0}}} &  \gate{H}\slice{\alert<3-5>{\QState{1}}} & \gate[wires=2][5em]{\mbox{$U_f$}}\gateinput{$x$}\gateoutput{$x$}\slice{\alert<6-8>{\QState{2}}} &\gate{H}\slice{\alert<9-10>{\QState{3}}} & \meter{\alt<30>{\alert{1}}{\alt<25>{\alert{0}}{0/1}}} \\
\lstick{\QZero{}} & \gate{X} &   \gate{H}   &  \qw\gateinput{$y$}\gateoutput{\Xor{y}{f(x)}} & \qw & \qw
\end{quantikz}}%
\end{center}
}{%
\only<1-7>{%
    \begin{align*}
        \visible<2->{\QState{0} &= \ket{01} \\}
        \only<3-5>{
        \visible<3->{\QState{1} &= \ket{+-} \visible<4->{= \frac{1}{2}\DQB{1}{-1}{1}{-1}}} \\}
        \visible<5->{\only<6->{\QState{1}}&= \frac{\ket{00}-\ket{01}+\ket{10}-\ket{11}}{2}}
    \end{align*}}%
\only<8-12>{%
\begin{align*}
    U_{f}(\ket{\textcolor{\RCone}{x}\textcolor{\RCtwo}{y}}) & = \ket{\textcolor{\RCone}{x}\,\Xor{\textcolor{\RCtwo}{y}}{f(\textcolor{\RCone}{x})}} \\
    \only<9>{U_{f}(\ket{\textcolor{\RCone}{0}\textcolor{\RCtwo}{0}}) & = \ket{\textcolor{\RCone}{0}\,\Xor{\textcolor{\RCtwo}{0}}{f(\textcolor{\RCone}{0})}}}
    \only<10>{U_{f}(\ket{\textcolor{\RCone}{0}\textcolor{\RCtwo}{1}}) & = \ket{\textcolor{\RCone}{0}\,\Xor{\textcolor{\RCtwo}{1}}{f(\textcolor{\RCone}{0})}}}
    \only<11>{U_{f}(\ket{\textcolor{\RCone}{1}\textcolor{\RCtwo}{0}}) & = \ket{\textcolor{\RCone}{1}\,\Xor{\textcolor{\RCtwo}{0}}{f(\textcolor{\RCone}{1})}}}
    \only<12>{U_{f}(\ket{\textcolor{\RCone}{1}\textcolor{\RCtwo}{1}}) & = \ket{\textcolor{\RCone}{1}\,\Xor{\textcolor{\RCtwo}{1}}{f(\textcolor{\RCone}{1})}}}
\end{align*}}
\only<13-18>{%
\begin{align*}
    \visible<13-15>{\Xor{0}{f(a)} &= f(a)} \\
    \visible<16->{\Xor{1}{f(a)} &= \Not{f(a)}}
\end{align*}}%
\only<22,27>{%
\BigSkip{}
\alert{Recall \[\ket{a\,b} = \ket{a}\ket{b} \]}
}%
\only<23-24>{%
\[ \alert{\ket{+} = \QPlus{}} \]
}%
\only<28-30>{%
\[ \alert{\ket{-} = \QMinus{}} \]
}%
\only<31>{%
\begin{center}
    \begin{tabular}{ccc}
        & \QState{2} & \QState{3} \\ 
        $f(x)$ & \multicolumn{2}{c}{(top qubit)}\\ \hline
    constant & \ket{+} & \QZero{} \\
    balanced & \ket{-} & \QOne{}
    \end{tabular}
\end{center}
}
}
\only<6-30>{
\Vskip{-2em}\begin{align*}
    \only<6-7>{\QState{2} &= U_f(\frac{1}{2}\left[\ket{00}-\ket{01}+\ket{10}-\ket{11}\right]) \\[1em]}
     \only<7->{\visible<8->{\QState{2}}   &= \frac{%
        \alt<14->{\ket{0\,f(0)}}{\alt<9->{\ket{0\,\alert<13>{\Xor{0}{f(0)}}}}{U_f(\ket{00})}}-\alt<17->{\ket{0\,\Not{f(0)}}}{\alt<10->{\ket{0\,\alert<16>{\Xor{1}{f(0)}}}}{U_f(\ket{01})}}  + \alt<15->{\ket{1\,f(1)}}{\alt<11->{\ket{1\,\alert<13-14>{\Xor{0}{f(1)}}}}{U_f(\ket{10})}}-\alt<18->{\ket{1\,\Not{f(1)}}}{\alt<12->{\ket{1\,\alert<16-17>{\Xor{1}{f(1)}}}}{U_f(\ket{11})}}}{2}}
\end{align*}}
\only<18-20>{%
\begin{itemize}
    \item<18-> For any $f(x)$, \QState{2} has 4 distinct terms: \ket{00}, \ket{01}, \ket{10}, and \ket{11}.
    \item<19-> $U_f$ determines which terms are positive and which are negative.
    \item<20-> If only the numerator contained \ket{0\,f(0)} - \ket{0\,f(1)}, then those terms would cancel if $f(0)$ is constant.
\end{itemize}
}
\only<21-25>{%
\Vskip{-2em}\TwoUnequalColumns{0.3\textwidth}{0.7\textwidth}{%
Suppose $f(0)=f(1)$.  \only<23-24>{%
Tensor-factored, this state tells us the \textcolor{\RCone}{top} and \textcolor{\RCtwo}{bottom} qubits' states.
}
\only<25>{Then the top qubit at \QState{2} is \ket{+}.  Measuring at \QState{3} yields \alert{\ket{0}}.}
}{%
\Vskip{-4em}\begin{align*}
    2\times\QState{2} &= \ket{0\,f(0)} - \ket{0\,\Not{f(0)}} + \ket{1\,\alert<21>{f(0)}} - \ket{1\,\Not{\alert<21>{f(0)}}} \\
  \visible<22->{  & = \TensProd{(\ket{0}+\ket{1})}{\ket{f(0)}} - \TensProd{(\ket{0}+\ket{1})}{\ket{\Not{f(0)}}} \\}
  \visible<23->{  & = \TensProd{\textcolor<23-24>{\RCone}{(\ket{0}+\ket{1})}}{\textcolor<23-24>{\RCtwo}{(\ket{f(0)}-\ket{\Not{f(0)}})}} \\}
  \visible<24->{ \QState{2} &= \TensProd{\textcolor<24>{\RCone}{\ket{+}}}{\textcolor<24>{\RCtwo}{(\ket{f(0)}-\ket{\Not{f(0)}})\frac{1}{\sqrt{2}}}} }
\end{align*}
}
}
\only<26-30>{%
\Vskip{-2em}\TwoUnequalColumns{0.3\textwidth}{0.7\textwidth}{%
Suppose $f(0) = \Not{f(1)}$. 
 \only<28-29>{%
Tensor-factored, this state tells us the \textcolor{\RCone}{top} and \textcolor{\RCtwo}{bottom} qubits' states.
}
\only<30>{Then the top qubit at \QState{2} is \ket{-}.  Measuring at \QState{3} yields \alert{\ket{1}}.}
}{%
\Vskip{-4em}\begin{align*}
    2\times\QState{2} &= \ket{0\,f(0)} - \ket{0\,\alert<26>{f(1)}} + \ket{1\,f(0)} - \ket{1\,\alert<26>{f(0)}} \\
  \visible<27->{  & = \TensProd{(\ket{0}-\ket{1})}{\ket{f(0)}} - \TensProd{(\ket{0}-\ket{1})}{\ket{f(1)}} \\}
  \visible<28->{  & = \TensProd{\textcolor<28-29>{\RCone}{(\ket{0}-\ket{1})}}{\textcolor<28-29>{\RCtwo}{(\ket{f(0)}-\ket{f(1)})}} \\}
  \visible<29->{ \QState{2} &= \TensProd{\textcolor<29>{\RCone}{\ket{-}}}{\textcolor<29>{\RCtwo}{(\ket{f(0)}-\ket{f(1)})\frac{1}{\sqrt{2}}}} }
\end{align*}
}
}
\only<31->{%
\MedSkip{}
\begin{itemize}
    \item<31-> We can determine if $f(x)$ is constant or balanced with a single query of $U_f$.  Classically two queries were needed.
    \item<32-> We perform a \href{https://physics.stackexchange.com/questions/184242/partial-measurement-and-the-math-behind-it}{partial measurement} of the system, measuring only the top qubit.
    \item<33-> This quantum circuit never created any entanglement:  each state could be tensor factored into the product of single qubit states.
    \item<34-> The effect of sending $y=\ket{-}$ into $U_f$ is manifest primarily on the top qubit.   We call this \emph{phase kickback} and study it next.
\end{itemize}
}

\end{frame}

\begin{frame}{What about the bottom qubit?}{Homework problem but let's take a quick look.}

\Vskip{-4em}\begin{align*} 
\invisible<5>{f(0) = f(1) \longrightarrow &  \QState{2} = \TensProd{\textcolor{\RCone}{\ket{+}}}{\textcolor{\RCtwo}{(\ket{f(0)}-\ket{\Not{f(0)}})\frac{1}{\sqrt{2}}}} \\}
\invisible<4>{f(0) \not= f(1) \longrightarrow & \QState{2} = \TensProd{\textcolor{\RCone}{\ket{-}}}{\textcolor{\RCtwo}{(\ket{f(0)}-\ket{f(1)})\frac{1}{\sqrt{2}}}}}
\end{align*}
\Vskip{-2em}\begin{itemize}[<+->]
  \item Tensor factored, we can see the values present on the \textcolor{\RCone}{top} and \textcolor{\RCtwo}{bottom} qubits.
  \item You will analyze the \textcolor{\RCtwo}{bottom qubit's} value for each of the above cases.
  \item We note here that the presence of \ket{-} as input for $y$ on the bottom qubit has the effect of changing phase on the \textcolor{\RCone}{\emph{top} qubit's} value on output:
  \begin{description}
     \item[$f(0)=f(1)$] $\textcolor{\RCone}{\ket{+}}=\TwoSupOp{\ket{0}}{\alert{\ket{1}}}{\alert{+}}$
     \item[$f(0)\not=f(1)$] $\textcolor{\RCone}{\ket{-}}=\TwoSupOp{\ket{0}}{\alert{\ket{1}}}{\alert{-}}$
  \end{description}
  \item This is called \emph{phase kickback}.
\end{itemize}
\end{frame}
\end{document}
